\documentclass[a4paper]{article}     
\usepackage{amsfonts}
\usepackage{amsmath}
\usepackage{amssymb}
\usepackage{graphicx}

\begin{document}

\noindent \large Auteur: Adriaan De Kaey \\
\noindent \large Studierichting: Informatica\\
\noindent \large Oefeningenreeks 1A nr. 4.6\\

\medskip

\normalsize

\textbf{Oefening 4.6}

Gegeven\\

$A=\left\{x\in\mathbb{R}^+ : \ln x = 1/n \text{ voor zekere }n\in\mathbb{N}_0\right\}$.\\

$\ln x = 1/n \Rightarrow x = e^{1/n}$.\\

Dus $A=\left\{e^{1/n} : n\in\mathbb{N}_0\right\}$.\\

Voor $n\ge 1$ geldt $1<e^{1/n}\le e$ en $e^{1/n}\to 1$.\\

$M\left(A\right)=\left[e,+\infty\right[$.\\

Dus $\max A = e$ en $\sup A = e$.\\

$m\left(A\right)=\left]-\infty,1\right]$.\\

Dus $\inf A = 1$ en er bestaat geen minimum.\\

$\inf A = 1$ en geen minimum.\\

\begin{thebibliography}{999}
%\bibitem{a} Referentie
\end{thebibliography}
\end{document} 
